\documentclass[10pt,a4paper]{amsart}
\usepackage[margin=19mm]{geometry}
\usepackage{amsmath,amssymb,mathtools}
\usepackage[scr=rsfs,cal=euler]{mathalfa}
\usepackage{xfrac}
\usepackage[unicode,hidelinks,bookmarks=false]{hyperref}
\newcommand{\ie}{i.e.\ }
\newcommand{\cf}{cf.\ }
\newcommand{\resp}{respectively\ }
\newcommand{\Subsection}{\subsection}
\providecommand{\nobreakdash}{\nobreak\hbox{-}}
\setlength{\emergencystretch}{2em}
\multlinegap0pt
\usepackage{bm}
\usepackage{bbm}
\usepackage{dsfont}
\usepackage{xspace}
\usepackage{cancel}
\usepackage{mathtools}
\usepackage{amsthm}
\makeatletter
\@ifundefined{claim}{\newtheorem{claim}{Claim}}{}
\makeatother
\makeatletter
\DeclareMathOperator{\N}{\textup{\textsf{N}}} 
\expandafter\ifx\csname bulletmidsentence\endcsname\relax
\newenvironment{bulletmidsentence}{\begin{itemize}[topsep=0mm, leftmargin=8mm, itemsep=1mm]
}{\end{itemize}}
\fi
\expandafter\ifx\csname rmkenumerate\endcsname\relax
\newenvironment{rmkenumerate}{\begin{enumerate}[align=left, leftmargin = 0mm, listparindent=0mm, labelwidth=2mm, itemindent=!, itemsep=1.5mm, labelsep=3mm, label=\textup{(\arabic*)}]
}{\end{enumerate}}
\fi
\expandafter\ifx\csname thmenumerate\endcsname\relax
\newenvironment{thmenumerate}{\begin{enumerate}[topsep=0mm, leftmargin=14mm,itemsep=1mm,label=\textup{(\roman*)}]
}{\end{enumerate}}
\fi
\makeatother
\title[Howson property and finitely generated intersection problem ]{Howson property and finitely generated intersection problem for monogenic inverse semigroups}
\author{Jung Won Cho, Craig Miller, Nik Ru\v{s}kuc}
\date{Source: arXiv:2607.19163v1 (CC BY 4.0)}
\begin{document}
\maketitle

\noindent\emph{Source and licence.} Howson property and finitely generated intersection problem for monogenic inverse semigroups. Version arXiv:2607.19163v1 (CC BY 4.0), distributed under the Creative Commons Attribution 4.0 International licence, as recorded in the arXiv licence metadata for this exact version and shown on the abstract page https://arxiv.org/abs/2607.19163v1. Full rights notice: licenses/arxiv-2607.19163-CC-BY-4.0.txt.\par\smallskip

\noindent\textit{Editorial scope.} Counted unit: the Claim whose source label is \texttt{cl:HowsonCounter1} in the article (source0.tex lines 676--679, source label \texttt{cl:HowsonCounter1}), delivered together with the article's own complete original proof of it (source0.tex lines 705--735, one \texttt{proof} environment of 31 lines). No mathematics is written, repaired or re-derived: statement and proof are the article's own text line for line. The proof is detached from the statement in the source: the article states the result at lines 676--679 and prints its proof at lines 705--735, and the proof environment's own header names the statement, so the association is the article's own. Required context retained uncounted: cl:HowsonCounter2 (lines 698--701). Every definition, display and label of those lines is retained unchanged. Omitted from this package and not redistributed: 1--675, 680--697, 702--704, 736--1759. Nothing inside a retained excerpt is omitted or altered: every retained body is delivered exactly as the article's lines are written, so no omission receipt of any kind accompanies this package. Citations: the retained excerpts carry no citation command at all, so no bibliography entry of the article is transcribed and no reference is redistributed. Reference adaptation: none was needed and none was made; every reference inside the delivered unit resolves inside it, so no unresolved reference or citation is produced. Printed slips kept literal: no printed word, symbol, hypothesis or number of the counted statement or of its proof was altered. Layout and class: the article's own class is \texttt{amsart} with packages including geometry, pdfsync, xcolor, amsmath, amsfonts, amssymb, amsthm, mathrsfs, hyperref, theoremref, tikz, float, inputenc, array; the wrapper is the lane's pinned \texttt{amsart} editorial wrapper at 10pt on A4, which restates the theorem-like environments the retained text needs and adds the article's own macro definitions that the retained text needs (\texttt{\textbackslash N}), with extra wrapper packages (bm, bbm, dsfont, xspace, cancel, mathtools). The delivered numbering is the unit's own. Assets: no retained span contains a figure environment, a graphics-inclusion command, a PDF-page inclusion, a table or a TikZ picture; no image file is copied into this package, the package declares no asset dependency and no image file is redistributed with it. Licence: version arXiv:2607.19163v1 of this article is distributed under the Creative Commons Attribution 4.0 International licence, as recorded in the arXiv licence metadata for this exact version and shown on the abstract page https://arxiv.org/abs/2607.19163v1. A full rights notice is delivered at licenses/arxiv-2607.19163-CC-BY-4.0.txt. Characters: every retained excerpt is pure ASCII, so the delivered engine, XeLaTeX with its default Latin Modern text font, reports no missing character.
\begin{claim}
\label{cl:HowsonCounter2}
Let $u\in E(T)$ be such that $u = x_1\dots x_n$ where each $x_i\in\{(-2, -2, 1), (0, 2, 2)\}$. Suppose that no proper initial subword of the product is an idempotent. Then, $u\rho$ is odd.
\end{claim}

\begin{claim}
\label{cl:HowsonCounter1}
Let $u\in \N(T)$ be such that $u = x_1\dots x_n$ where each $x_i\in\{(-2, -2, 1), (0, 2, 2)\}$. Suppose that no proper initial subword of the product is an idempotent. Then, $u\rho = 1$.
\end{claim}

\begin{proof}
The assumption means that in all proper initial subwords the terminal vertex is on the same side in 
relation to the initial vertex. We split our considerations into two cases depending on the value of $x_1$.

\emph{Case 1: $x_1 = (-2, -2, 1)$.} Note that every proper initial subword must be negative and that $x_n = (0, 2, 2)$. The initial subword $x_1\dots x_{n-1}$ satisfies the assumptions from Claim~\ref{cl:HowsonCounter1}. Notice that $(x_1\dots x_{n-1})\mu = -2$, and then the final multiplication by $(0, 2, 2)$ does not increase the value of $\rho$. Hence, $u\rho = (x_1\dots x_n)\rho = (x_1\dots x_{n-1})\rho = 1$ by Claim~\ref{cl:HowsonCounter1}.

\emph{Case 2: $x_1 = (0, 2, 2)$.} In this case, every proper initial subword is positive and $x_n = (-2, -2, 1)$. Let 
\[
k := \min\{1\leq i \leq n: (x_1\dots x_i)\rho = u\rho\}.
\]
We claim that $x_k$ must be $(-2, -2, 1)$. Suppose otherwise. Then, 
\[
\bigl((x_1\dots x_{k-1})(0, 2, 2)\bigr)\rho = \max\bigl((x_1\dots x_{k-1})\rho, (x_1\dots x_{k-1})\mu + 2\bigr) = (x_1\dots x_{k-1})\mu + 2,
\]
by minimality of $k$; in particular, $(x_1\dots x_k)\rho = (x_1\dots x_k)\mu$. Since $x_k$ is not the final letter, $x_1\dots x_k$ is followed by $(0, 2, 2)$ or $(-2, -2 ,1)$. 
In each case we have:
\begin{align*}
    &\bigl((x_1\dots x_k)(0, 2, 2)\bigr)\rho = \max\bigl((x_1\dots x_k)\rho, (x_1\dots x_k)\mu + 2\bigr) = (x_1\dots x_k)\mu + 2 > (x_1\dots x_k)\rho;\\
    &\bigl((x_1\dots x_k)(-2, -2, 1)\bigr)\rho = \max\bigl((x_1\dots x_k)\rho, (x_1\dots x_k)\mu+1)\bigr) = (x_1\dots x_k)\mu+1>(x_1\dots x_k)\rho.
\end{align*}
Therefore
\[
(x_1\dots x_n)\rho\geq (x_1\dots x_{k+1})\rho> (x_1\dots x_k)\rho,
\]
which contradicts the choice of $k$, and thus establishes the claim that $x_k = (-2, -2, 1)$.  
The following is now straightforward:
\[
u\rho = \bigl((x_1\dots x_{k-1})(-2, -2, 1)\bigr)\rho = \max\bigl((x_1\dots x_{k-1})\rho, (x_1\dots x_{k-1})\mu+1\bigr) = (x_1\dots x_{k-1})\mu + 1,
\]
which is odd since $T\mu = 2\mathbb{Z}$. This completes the proof of Claim \ref{cl:HowsonCounter2}.
\end{proof}

\end{document}
